\subsection{Source-bound identifications with known constructions}
The new positive inventory includes tensor products of previously
constructed unitary categories and representation categories of finite
groups. Each entry is checked by a unit-preserving basis bijection against
all $8^3=512$ source coefficients, including noncommutative products.
The certificate specifies the direction of the basis map. Equality of
dimension vectors alone is never used as a model identification.

For the Frobenius-group entry $\sha{R8_3_342}$, the finite orbit-convolution
calculation constructs $\operatorname{Rep}(C_{13}\rtimes C_6)$ with its
faithful action. The $\operatorname{PSL}_2(11)$ entry
$\sha{R8_3_724}$ is reconstructed from its exact finite-group class algebra.
The $E_8$ module-category entry $\sha{R8_3_3278}$ uses the standard unitary
$E_8$ algebra and its module fusion rules \cite{KO}. The remaining existing
near-group, Fibonacci and $E_8$ product entries use the unitary
$\operatorname{NG}(C_3,3)$ construction
\cite[Proposition~A.4.1]{CMSIzumiXu} or the adjoint of the unitary
$E_8$ construction \cite{KO}, followed by the indicated Deligne product.
Their complete tensor comparisons are in the model catalogue.
These comparisons establish the based-ring identification;
categorical existence uses the cited construction and its stated scope.

\subsection{Unitary scope of the Izumi--Xu extension}
The fusion rules of $\sha{R8_3_448}$ are those of \cite[Appendix~A.6.1]{CMSIzumiXu}; in particular $Vg=g^2V$ is essential. The unitary construction is supplied by the symmetrically self-dual $2221$ planar algebra of \cite[Section~5.1, Theorem~5.5]{MorrisonPenneysSpokes}. The extension classification in \cite[Theorem~4.9, Remarks~4.10--4.11]{GJSCyclic} identifies the resulting graded extension with the displayed fusion rules: the two associator choices have the same tensor. Thus the earlier conservative pseudo-unitary label is upgraded to unitary existence.

\subsection{Two equivariantizations of a near-group category}
Use the normalized unitary $C_3+6$ solution of \cite[Example~13.10]{IzumiNear}, with $(x,y)=(\sqrt{\sqrt3/24},0)$. Lemma~13.1 describes actual automorphisms $\gamma_{\omega,\theta,u}$, with $(\omega,\theta,u)$ and $(-\omega,\theta,-u)$ identical. The paragraph before Lemma~13.2 gives the strict conjugation tensor functor. Take $\theta=\mathrm{id}$ and either $(\omega,u)=(1,\mathrm{diag}(1,-1))$ or $(i,R_{\pi/2})$. Both have order two; the latter squares to $\gamma_{-1,\mathrm{id},-I}=1$. Their actions on the six-dimensional self-fusion space have three eigenvalues of each sign. Their actions on the maps $S_g$ have signs $+1$ and $-1$, respectively. Hence the complete rules are
\[
 (g,e)(h,f)=(g+h,e+f),\quad (g,e)X_f=X_f(g,e)=X_{e+f},
 \qquad X_eX_f=\sum_{g\in C_3}(g,e+f+\epsilon)+3X_0+3X_1.
\]
The cases $\epsilon=0,1$ identify $R8\_3\_334,R8\_3\_274$ by all512 coefficients. These are coherent unitary actions, not abstract based-ring automorphisms.

\subsection{Cyclic permutation of three Fibonacci factors}
The canonical $C_3$ action on $\mathrm{Fib}^{\boxtimes3}$ has two fixed simple triples, each splitting into three equivariant simples, and two free orbits. For a fixed output, the multiplicity is the character average on its fusion space; a free input orbit has zero nonidentity trace. For a free output the coefficient is the ordinary multiplicity at a representative. These formulas reconstruct all512 coefficients of $\sha{R8_3_475}$ and give its unitary realization.

\subsection{A Morita dual of the nonsplit $2D2$ extension}
In the known nonsplit extension $\sha{R8_2_37}$, take the Q-system $A=1+r$ from \cite[Lemma~4.1]{Grossman}. The neutral component of ${}_A\mathcal E_A$ is $TT_3$, with basis $1,p,q,pq,qp,pqp$ and relations $p^2=1+p$, $q^2=1+q$, $pqp=qpq$; see \cite[Definition~2.5]{EdieGolden}. The odd free-bimodule Gram matrix, in the order $(a,a^3,ar,a^3r)$, is
\[
 \begin{pmatrix}1&1&3&3\\1&1&3&3\\3&3&13&13\\3&3&13&13\end{pmatrix}.
\]
Its nonnegative integral factorizations have the column $(1,1,3,3)^T$ and either $(0,0,2,2)^T$ or four copies of $(0,0,1,1)^T$. After dividing underlying dimensions by $\operatorname{FPdim}(A)$, the latter would give four odd simples of dimension $\varphi$, but multiplication by $p$ requires the unavailable dimension $\varphi^2$. The two remaining odd dimensions are therefore $2\varphi,2\varphi^2$. Left and right multiplication by either $p$ or $q$ is forced to be $\left(\begin{smallmatrix}0&1\\1&1\end{smallmatrix}\right)$. Reciprocity determines odd--odd products and yields precisely $\sha{R8_3_715}$.

\subsection{The conformal $G_2$ algebra Morita dual}
Let $\mathcal C=\mathcal C(G_2,3)$ and let $A$ be the conformal algebra for $G_{2,3}\subset E_{6,1}$. The category ${}_A\mathcal C_A$ is unitary by \cite[Theorems~2.21,2.28, Corollary~3.32]{GuiQSystems}. Write $d=(3+\sqrt{21})/2$. The two chiral Izumi--Xu subcategories have pointed local intersection $C_3$, by \cite[Appendix~A.4]{CMSIzumiXu} and \cite[Section~3.3, equation~(16), Definition~3.12, Remark~3.13]{DMNO}. Their noninvertibles $X,Y$ satisfy $X^2=1+g+g^2+3X$ and $Y^2=1+g+g^2+3Y$. Rigidity gives $(XY,XY)=3$; dimension equality forces three further simples $Z_i$, each of dimension $d+1$.

The $g$-action on this triple is nontrivial: otherwise the symmetric integral multiplicity block $B$ has row sums three and satisfies $B^2+J=3I+3B$. Its trace is six, so its off-diagonal sum would be three, a parity contradiction. Write $Z_i=g^iZ_0$ with $Z_0$ self-dual. Associativity gives $XY=YX=\sum_i Z_i$, $XZ_i=Z_iX=Y+\sum_j Z_j$, and $YZ_i=Z_iY=X+\sum_j Z_j$. The commutative orientation is impossible: the modular center has only self-dual simples, and \cite[Lemma~8.49]{ENO} would force all simples of the commutative based ring to be self-dual, contrary to $g^*=g^2$. Hence $Z_ig=g^{-1}Z_i$. Reciprocity reduces the remaining products to ten nonnegative triples with sum three; exact associativity leaves only $(1,1,1)$, so
\[
 Z_iZ_j=g^{i-j}+X+Y+\sum_k Z_k.
\]
The reconstructed full tensor is $\sha{R8_3_458}$.

\subsection{A twisted group-theoretical construction}
% Add bibliography key CartanEilenberg: H. Cartan and S. Eilenberg,
% Homological Algebra, Princeton University Press (1956), Chapter XII, Section 10.
% GelakiNaidu is already present in the rank-eight bibliography.
\begin{theorem}[A twisted $A_6/A_5$ model]\label{thm:m3-twisted-a6}
There are $\mu_2$-valued cochains $\omega\in Z^3(A_6,\mu_2)$ and
$\psi\in C^2(A_5,\mu_2)$ with $d\psi=\omega|_{A_5}$ such that the
unitary group-theoretical category
$\mathcal C(A_6,\omega,A_5,\psi)$ realizes $\sha{R8_3_3342}$.
Its simple dimensions are $1,3,3,4,5,10,10,10$.
\end{theorem}
\begin{proof}
We give an exact finite certificate for the relative cocycle and then
derive the fusion rules. Let $P=\langle r,s\rangle\cong D_8$ in $A_6$,
on letters $0,\ldots,5$, with
\[
 r=(0123)(45),\qquad s=(02)(45).
\]
Write $r^is^j=(i,j)$, and set
$E=\{(i,j):i\equiv j\pmod2\}$,
$K=\{(i,j):i\equiv0\pmod2\}$, and $t=s$.
Then $E$ fixes both $4$ and $5$, while $K$ is the other Klein four
subgroup. On $K$, put $x(i,j)=i/2$, $y(i,j)=j$, with values in
$\mathbb F_2$. The normalized bar equations
\[
 dw=0,\qquad w|_E=0,\qquad w|_K=x^2y+xy^2+dv
\]
have an explicitly saved solution: there are $352$ variables,
$2455$ equations, and rank $312$.

For each of the $360$ elements $g\in A_6$, the certificate supplies a
normalized $2$-cochain witnessing equality, modulo coboundaries, of
$w$ and its $g$-conjugate on $P\cap gPg^{-1}$. The intersection orders
$1,2,4,8$ occur $256,64,32,8$ times, respectively. Thus the
Cartan--Eilenberg stable-elements theorem
\cite[Chapter XII, Section 10]{CartanEilenberg} supplies a class
$c\in H^3(A_6,\mathbb F_2)$ restricting to $[w]$.
For $H=A_5$ fixing $5$, $E$ is a Sylow $2$-subgroup of odd index $15$;
restriction and transfer give injectivity of $H^3(H,\mathbb F_2)$
into $H^3(E,\mathbb F_2)$. Hence $c|_H=0$.
Choose a normalized trivializing $2$-cochain on $H$, extend it to
$A_6$, and gauge the global cocycle to vanish identically on $H$.
Exponentiation by $a\mapsto(-1)^a$ gives $\omega|_H=1$, so we may
take $\psi=1$.

The off-diagonal stabilizer is $L=H\cap tHt^{-1}=A_4$.
The multiplier convention of \cite[Section 5.1]{GelakiNaidu}, restricted
to the involutions $a,b\in E$, has additive form
\[
 \alpha(a,b)=w(a,b,t)+w(a,bt,t^{-1}bt)
                   +w(abt,t^{-1}bt,t^{-1}at).
\]
For $u=r^2$ and $v=rs$, the exact affine certificate proves
$\alpha(u,v)+\alpha(v,u)=1$: it gives a sum of original equations
equal to this functional plus $1$. The identity therefore holds for
every solution of the displayed affine system, including the
restriction of the globally gauged cocycle. This last restriction
need only be cohomologous to the saved local representative.
Consequently the multiplier has commutator $-1$ on a commuting pair,
so its class in $H^2(A_4,\mathbb C^\times)$ is nontrivial.
This relative computation is essential: the nonzero Dickson cubic
itself maps to zero in $H^3(K,\mathbb C^\times)$.

There are two double cosets. The diagonal sector is
$\operatorname{Rep}(A_5)$. The other sector has the three spin
representations of the binary tetrahedral group $2T$:
\[
 \rho_i=\rho_0\otimes\chi^i,\qquad i\in\mathbb Z/3,
 \qquad\det\rho_i=\chi^{2i},
\]
where $\rho_0$ is the defining $SU(2)$ representation and $\chi$
generates the character group of $A_4$. The nontrivial $\mu_2$
extension is unique: $H^2(A_4,\mathbb F_2)$ is the one-dimensional
$C_3$-invariant part of $H^2(V_4,\mathbb F_2)$.
The three spin representations have squared dimensions summing to
$12$, so they exhaust the projective simples. Their corresponding
objects $X_i$ have dimension $[A_5:A_4]\dim\rho_i=10$.

Every $X_i$ is self-dual. Reversing the endpoints conjugates $A_4$
by an odd permutation and therefore inverts its character group.
Representation duality also inverts determinants, so the two
inversions cancel. All remaining transport scalars lie in $\mu_2$
and hence do not change two-dimensional determinants. The three
distinct determinants distinguish the $X_i$. The five diagonal
simples are also self-dual; hence all fusion coefficients are
symmetric in their three indices.

At distinct endpoints, composing two off-diagonal kernels sums over
four intermediate points, a transitive $A_4$-set with stabilizer
$C_3$. Thus the off-sector product is governed by induction from
this $C_3$. There is a unique $\mu_2$ trivialization of a multiplier
on $C_3$, since $H^1(C_3,\mu_2)=H^2(C_3,\mu_2)=0$.
The order-three lift in $2T$ has eigenvalues $\zeta,\zeta^2$ on
$\rho_0$, so $\rho_0|_{C_3}=\chi+\chi^2$. This is preserved by
either automorphism of $C_3$; a residual $\mu_2$ transport character
is necessarily trivial. Equivalently, transport preserves the
trivial determinant, whereas a nontrivial $C_3$ character twist
would change it. Therefore
\[
 N_{X_0X_0}^{X_0}
 =\langle\chi+\chi^2,(\chi+\chi^2)^2\rangle_{C_3}=2.
\]
For any pair $i,j$, the tensor of their two-character restrictions
has multiplicities $(2,1,1)$ in some order. Pairing with the three
spin restrictions gives off-sector multiplicities $(2,3,3)$ in some
order. Symmetry and the anchor $N_{X_0X_0}^{X_0}=2$ have a unique
completion (enumerated among the $3^6$ choices for unordered pairs):
\[
 N_{X_iX_j}^{X_k}=3-\delta_{i+j+k,0\pmod3}.
\]

The restrictions of the five diagonal simples to $A_4$ are
$1,3,3,1+3,\chi+\chi^2+3$. As
$3\otimes\rho_i=\rho_0+\rho_1+\rho_2$, their actions on the
three off-sector objects are $I,J,J,I+J,2J-I$, where $J$ is the
all-ones matrix. Reciprocity supplies the diagonal coefficients
of off/off products. The remaining tensor is the ordinary $A_5$
character tensor. This determines every coefficient.

The verifier constructs that tensor independently, checks the
dimension and associativity equations, and matches all $512$
coefficients to the authoritative source under the basis map
$[0,1,2,6,7,3,4,5]$ from model to source, whose source-to-model inverse
is $[0,1,2,5,6,7,3,4]$. All source hashes are recomputed. Finally,
unit-modulus cocycles make the pointed category and its subgroup
bimodule construction unitary.
\end{proof}

\begin{remark}
The comparison tensor $\sha{R8_3_3344}$ has
$N_{XX}^{X}=3$ for each dimension-ten simple. Thus it is not the
ring of this particular construction. This comparison is not a
general obstruction to other constructions of that tensor.
\end{remark}

\subsection{A Q-system realization of \texorpdfstring{$\sha{R8_3_721}$}{R8-3-721}}
\label{sec:r721-qsystem-model}
The earlier center calculations left this ring undecided. We now construct
an actual unitary fusion category and identify its Grothendieck ring within
the supplied complete census. The existence argument uses an inherited
Q-system; agreement of modular data is not used to manufacture an algebra.

\begin{theorem}[A cyclic generalized-Haagerup model]
\label{thm:r721-qsystem-model}
Let $\mathcal C$ be the $C_2$ de-equivariantization of the unitary
$C_8$ generalized-Haagerup category in
\cite[Example~4.4 and Section~6]{GrossmanIzumiCentersNext}.
Let $A=\one\oplus\rho$ be the image of its Q-system from
\cite[equations~(2.21)--(2.22)]{GrossmanIzumiCentersNext}.
Then $\mathcal D={}_A\mathcal C_A$ is a unitary fusion category of rank eight,
with trivial group of invertibles and FP dimensions
\[
 1,a,a,a+1,a+1,2a+1,2a,2a+2,
 \qquad a=\frac{3+\sqrt{17}}2.
\]
All its fusion coefficients are at most three. The exact completion of the mandatory induction equations below identifies
$K_0(\mathcal D)\cong\sha{R8_3_721}$.
\end{theorem}

We separate the categorical construction, the finite coefficient bound,
and the exact fusion completion. The census is used only at the end to
name the uniquely generated based-ring class.

\begin{lemma}[An inherited-vertex obstruction]
\label{lem:cyclic-gh-invertibles-next}
Suppose an existing unitary generalized-Haagerup category for $C_{4m}$
satisfies the de-equivariantization hypotheses of
\cite[Section~6]{GrossmanIzumiCentersNext}, including
$\epsilon_{2m}(h)=(-1)^h$, and has the Q-system $\one\oplus\rho$
whose nonzero $\rho\to\rho^2$ vertex is proportional to $t_0$.
In its $C_2$ de-equivariantization, the Morita dual of this inherited algebra
has no nontrivial invertible object.
\end{lemma}
\begin{proof}
Write $g$ for the generator of the quotient group $C_{2m}$ and
$d=2m+\sqrt{4m^2+1}$. An invertible bimodule has underlying FP dimension
$1+d$. Since the source simple dimensions are $1$ and the irrational number
$d$, its underlying object has one invertible and one noninvertible
summand. If its invertible summand is $g^j$, free-module adjunction and
 equality of dimensions identify its underlying left and right modules with
$A g^j$ and $g^j A$, respectively. Their noninvertible summands are
$g^{-j}\rho$ and $g^j\rho$, so $2j=0$ in $C_{2m}$.

For $j=0$, the bimodule is an automorphism-twisted copy of $A$.
A unit-preserving algebra automorphism acts by a scalar $z$ on $\rho$.
The nonzero $\rho^2\to\one$ and $\rho^2\to\rho$ multiplication components
impose $z^2=1$ and $z^2=z$, hence $z=1$.

For $j=m$, use the crossed-product realization of
\cite[Section~6]{GrossmanIzumiCentersNext}. Its implementing unitary
$\lambda$ satisfies
\[
 \lambda^2=1,\qquad \rho'(\lambda)=\lambda,\qquad
 \alpha'_{2m}=\operatorname{Ad}(\lambda),\qquad
 \alpha'_m(t_0)=t_{2m}.
\]
Conjugation by $\alpha'_m$ sends $\rho'$ to
$\operatorname{Ad}(\lambda)\rho'$. The intertwiner
$u=\lambda:\rho'\to\operatorname{Ad}(\lambda)\rho'$ has
$u\otimes u=1$. Transporting the conjugate algebra to the object
$\one\oplus\rho'$ changes its vertex line from $\C t_0$ to
$\C t_{2m}\lambda$. These lines are orthogonal by the Cuntz relations.
A unit-preserving object isomorphism acts scalarly on the simple
$\rho'$ summand and cannot identify the two vertex lines. Thus
$g^m A g^{-m}\not\cong A$ as algebras. An invertible bimodule with underlying
free module $g^m A$ would give precisely such an algebra isomorphism,
contradicting this calculation.
\end{proof}
This lemma is conditional on existence of the stated generalized-Haagerup
category and Q-system. It makes no existence assertion for arbitrary $m$.

\begin{proof}[Proof of Theorem~\ref{thm:r721-qsystem-model}]
In the case $m=2$, the source category and Q-system are supplied by the
cited construction. The strong tensor de-equivariantization functor
preserves the Q-system, and its image $\rho$ is simple. Hence
$A=\one\oplus\rho$ is connected and separable, and its bimodule category
exists and is unitary. The source fusion rules are
\[
 g^4=1,\qquad g\rho=\rho g^{-1},\qquad
 \rho^2=\one+2\sum_{i=0}^3g^i\rho,
 \qquad \FPdim(\rho)=4+\sqrt{17}.
\]
Morita equivalence gives
$\FPdim(\mathcal D)=136+32\sqrt{17}$.

For free bimodule induction $I(X)=A\otimes X\otimes A$, adjunction gives
\[
 K_{ij}=\dim\operatorname{Hom}_{\mathcal D}(I(X_i),I(X_j))
       =[A X_j A:X_i].
\]
Every simple bimodule occurs in some $I(X)$, since its action map from the
free bimodule on its underlying object is a split epimorphism.
The known tensor-unit row is $(1,0,0,0,1,0,0,0)$ in source order
$(1,g,g^2,g^3,\rho,g\rho,g^2\rho,g^3\rho)$.
The remaining Gram matrix has rank four. Its complete finite factorization
checks $40$ pivot vectors and $14256$ multiplicity choices; it yields seven
row types and exactly four factors, of ranks $6,8,9,11$.

The known center in
\cite[Theorem~6.8 and Example~6.10]{GrossmanIzumiCentersNext}, equivalently
\cite[Definition~6.1 and Section~6.1.6]{GrossmanIzumiMetricNext}, removes
ranks nine and eleven. Indeed its twist-one support consists of six objects,
and the canonical Lagrangian multiplicity vector has form
$(1,p,x,y,z,w)$ on that support. The equations $S\ell=\ell$ and
$T\ell=\ell$ force
\[
 p=1,\qquad x+y=2,\qquad z+w=2.
\]
All nine nonnegative integral choices satisfy the full $S$ equations, and
\[
 \rank(\mathcal D)=\langle\ell,\ell\rangle
       =2+x^2+y^2+z^2+w^2\in\{6,8,10\}.
\]
Thus only the rank-six and rank-eight Gram factors remain.
The rank-six factor has an extra invertible bimodule with underlying object
$g^2+g^2\rho=g^2A$. Lemma~\ref{lem:cyclic-gh-invertibles-next} excludes it.
The actual dual therefore has the unique rank-eight factor, and dividing
its underlying dimensions by $\FPdim(A)$ gives the dimensions in the theorem.

It remains to bound its unknown fusion coefficients, without assuming a
bound in the construction. Let $B_{i x}=[U(x):X_i]$ be the unique rank-eight
forgetful matrix. Free induction satisfies
\[
 I(X)\otimes_A I(Y)\cong I(X\otimes A\otimes Y).
\]
For every source pair $(i,j)$ and dual output $z$, this gives
\[
 \sum_{x,y}B_{i x}B_{j y}N_{xy}^{z}
   =\sum_{a\in\{1,\rho\}}\sum_{u,t}
        N^{\mathcal C}_{i a}{}^{u}
        N^{\mathcal C}_{u j}{}^{t} B_{t z}.
\]
Together with the unit, duality, Frobenius reciprocity and both rational
components of the dimension character, these are $832$ exact linear
equations before redundant rows are removed. There are four possible
dualities on the two equal-dimension pairs; they give respectively
$120,112,112,104$ coefficient orbits. All four are retained.

Write one reduced system as $E n=b$, where $n\ge0$ consists of those fusion
coefficients. For every coordinate $j$, the supplied rational covector
$y_j$ satisfies
\[
 y_j E\ge e_j^{\mathsf T},\qquad y_j b\le3.
\]
Therefore $n_j\le y_j E n=y_j b\le3$. The certificates check every inequality
in exact rational arithmetic for all four dualities; numerical optimization
is used only to propose the covectors. This proves the required uniform
multiplicity bound for the actual category.

Finally, we identify the actual category directly from its mandatory
linear equations, without using completeness of the supplied census.
After Frobenius and duality orbits are merged, the four possible dualities
on the two equal-dimension pairs have $120,112,112,104$ coefficient
variables. Impose the unit equations, both rational dimension coordinates
and all $512$ free-induction identities
\[
 \sum_{x,y}B_{ix}B_{jy}N_{xy}^z
 =\sum_{u,t}\sum_{a\in\{1,\rho\}}
        N^{\mathcal C}_{ia}{}^uN^{\mathcal C}_{uj}{}^tB_{tz}.
\]
The coefficient bounds above apply to every orbit in every duality branch.
Rational elimination leaves $22,19,19,16$ free coordinates, each itself
an actual fusion coefficient in $\{0,1,2,3\}$. Every other coefficient has
form $(c+\sum_j a_jx_j)/D$ with integral $c,a_j$ and $D>0$.
For each remaining value of $x_j$, interval arithmetic over the other
finite domains removes it only if the resulting affine interval cannot
meet $[0,3D]$; the necessary gcd congruence checks integrality. Repeating
these deductions and branching over every retained value is exhaustive.
There is no solver timeout, floating point or associativity pruning.
The complete search is
\[
\begin{array}{cc|rr}
 a\text{-pair duality}&b\text{-pair duality}&\text{nodes}&\text{solutions}\\\hline
 \text{fixed}&\text{fixed}&3&0\\
 \text{fixed}&\text{exchanged}&3&0\\
 \text{exchanged}&\text{fixed}&5&2\\
 \text{exchanged}&\text{exchanged}&3&0
\end{array}
\]
Thus the necessary linear equations alone leave two tensors. Both satisfy
all $4096$ associativity identities, all induction identities, dimensions,
unit and Frobenius reciprocity. They differ by exchanging $b_0,b_1$, so
form one based-ring class. Only now do we compare with the supplied census:
all $512$ coefficients match $\sha{R8_3_721}$ under the basis map
$(0,1,2,4,6,3,5,7)$, up to the harmless equal-dimension relabeling.
This proves the claimed identification independently of census completeness.
\end{proof}

\paragraph{Exact evidence and scope.}
The rank enumeration is
\path{evidence/rank8_next_20261021/semion_fermion/construction/r721_bimodule_rank.py}.
The coefficient system and rational certificates are in the same evidence
tree's \path{models/} directory:
\begin{quote}\small
\path{r721_dual_fusion.py}\quad and\quad
\path{r721_dual_multiplicity_bound.py}.
\end{quote}
The solver-free exact completion and complete branch traces are in
\begin{quote}\small\path{evidence/correctness_20261025/r721_exact_completion.py}.\end{quote}
An independent bounded-completion implementation agrees. The earlier
dimension crosswalk is corroboration only, not a completeness premise.
Normal and optimized-Python receipts, independent reconstructions, the
ordered tensor and the cited Q-system argument accompany these files.
The identification proves existence, and does not classify all associators
on this ring or assert that the Morita class is new.
